\chapter{Conclusão}


Com a realização deste trabalho no âmbito da unidade curricular de Laboratórios de Informática, foi possível adquirir e consolidar conhecimentos fundamentais relacionados com o controlo de versões, o trabalho colaborativo e a utilização de ferramentas essenciais no desenvolvimento de projetos informáticos.

\vspace*{0.5cm}

A utilização do GitHub revelou-se um elemento central para a organização e distribuição das tarefas entre os membros do grupo, permitindo uma gestão eficaz do projeto e um acompanhamento claro da evolução do trabalho. Através do histórico de commits, foi possível identificar as contribuições individuais, garantindo uma maior responsabilidade e colaboração entre todos os elementos do grupo.

\vspace*{0.5cm}

O recurso ao Visual Studio Code, em conjunto com uma máquina virtual Ubuntu, proporcionou uma experiência prática próxima da realidade profissional, permitindo a execução de commits diretamente através da consola. Esta abordagem contribuiu para o reforço das competências ao nível da linha de comandos, bem como para uma melhor compreensão do funcionamento do sistema de versionamento Git em ambientes Linux.

\vspace*{0.5cm}

De um modo geral, os objetivos propostos para este trabalho foram cumpridos com sucesso, tendo sido possível desenvolver não só competências técnicas, mas também capacidades de organização, comunicação e trabalho em equipa. A realização deste projeto constituiu, assim, uma experiência enriquecedora e relevante para a formação académica dos membros do grupo, preparando-os para desafios futuros na área da informática.